\chapter{基本的な積分と式の整理}\label{章:基本積分}

前章では，定積分を計算するために原始関数を探せばよいことが分かった。そこで今度は，目の前の関数に対して，その原始関数をどう見つけるかを考える。たとえば $x^2$ なら，微分すると $3x^2$ になる $x^3$ を思い出し，係数を直せばよい。ところが $(x+1)^2$ や $\frac{x^2}{x+1}$ は，見た形が違うだけで，同じ方法へ進めることがある。

この章では，まず微分公式を逆向きに読む。そのうえで，展開，割り算，平方完成，部分分数，三角恒等式を使い，知っている原始関数が見える形へ式を整える。新しい積分法を次々に覚える前に，式を変えるだけで何ができるかを確かめよう。

\section{微分公式を逆向きに読む}\label{節:基本公式の逆}

$F'(x)=f(x)$ ならば，$F$ は $f$ の原始関数である。この関係を使うとき，最初から正しい係数まで思い付く必要はない。まず形の近い候補を微分し，元の関数との差を直すのでよい。たとえば
\[
 \frac{\dd}{\dd x}x^5=5x^4
\]
だから，$x^4$ の原始関数を得るには，$x^5$ を $5$ で割る。
\[
 \int x^4\dd x=\frac{x^5}{5}+C.
\]
同じ理由で，$r\ne-1$ のとき
\[
 \int x^r\dd x=\frac{x^{r+1}}{r+1}+C
\]
となる。指数を一つ増やし，増やした指数で割るのは，微分で指数が一つ減り，その指数が係数として出ることを元に戻しているのである。

ここで $r=-1$ だけは別に考える必要がある。上の式へ $r=-1$ を代入すると分母が $0$ になり，意味を持たない。また，$x^0=1$ を微分しても $x^{-1}$ は現れない。代わりに，対数の微分
\[
 \frac{\dd}{\dd x}\log x=\frac1x\qquad(x>0)
\]
を使う。$x<0$ では $\log(-x)$ を微分すれば
\[
 \frac{\dd}{\dd x}\log(-x)
 =\frac1{-x}\,(-1)=\frac1x
\]
となるから，二つの形は $\log|x|$ とまとめられる。

\begin{theorem}{基本的な不定積分}
次の公式は，右辺が微分でき，左辺の関数が定義される区間で成り立つ。
\[
\begin{array}{c|c}
 f(x)&\int f(x)\dd x\\ \hline
 x^r\quad(r\ne-1)&\frac{x^{r+1}}{r+1}+C\\[5pt]
 \frac1x&\log|x|+C\\[5pt]
 \e^x&\e^x+C\\[3pt]
 a^x\quad(a>0,\ a\ne1)&\frac{a^x}{\log a}+C\\[5pt]
 \sin x&-\cos x+C\\[3pt]
 \cos x&\sin x+C\\[3pt]
 \frac1{\cos^2x}&\tan x+C
\end{array}
\]
実数 $r$ に対するべきの公式は，まず $x>0$ で考える。整数のべきなど，負の $x$ でも定義される場合は，その定義域内の区間にも適用できる。
\end{theorem}

この表は，右の式を微分して確かめられる。たとえば $a^x$ の微分では $\log a$ が掛かるので，積分では $\log a$ で割る。$a=1$ を除くのは，この割り算ができないためであり，$1^x=1$ の積分そのものは $x+C$ である。

不定積分の $C$ は，考えている一つの区間上での任意定数である。$\frac1x$ の定義域は $0$ の左右に分かれる。左側で加える定数と右側で加える定数が，同じである必要はない。図\ref{図:対数の区間}では，この区間の分かれ方を数直線に示した。原始関数を一つ選ぶ場合と，定義域全体のすべての原始関数を述べる場合とを区別しよう。

\begin{figure}[htbp]
 \centering
 \begin{tikzpicture}[x=1cm,y=1cm]
  \draw[-{Stealth}] (-4,0)--(4,0) node[right] {$x$};
  \draw[very thick,blue] (-3.7,0)--(-0.12,0);
  \draw[very thick,blue] (0.12,0)--(3.7,0);
  \filldraw[fill=white] (0,0) circle (2pt);
  \node[below] at (0,-0.08) {$0$};
  \node[above] at (-2,0.18) {$\log(-x)+C_-$};
  \node[above] at (2,0.18) {$\log x+C_+$};
  \node[below] at (-2,-0.35) {$x<0$};
  \node[below] at (2,-0.35) {$x>0$};
 \end{tikzpicture}
 \caption{$0$ を越えて一つの区間にはできない。定数は左右で独立に選べる。}
 \label{図:対数の区間}
\end{figure}

\begin{example}{根号と分母をべきの形へ直す}
$x>0$ において，$\int\left(2\sqrt{x}+\frac3{x^2}\right)\dd x$ を求めよ。
\end{example}

根号のままでは公式と見比べにくいので，$\sqrt{x}=x^{\frac12}$，$\frac1{x^2}=x^{-2}$ と書く。二つの項は別々に積分できるから，
\[
\begin{aligned}
 \int\left(2\sqrt{x}+\frac3{x^2}\right)\dd x
 &=2\int x^{\frac12}\dd x+3\int x^{-2}\dd x\\
 &=2\cdot\frac23x^{\frac32}
   +3\cdot\frac{x^{-1}}{-1}+C\\
 &=\frac43x^{\frac32}-\frac3x+C.
\end{aligned}
\]
分母にある式をいつも対数で積分するわけではない。$\frac1x$ は指数が $-1$ なので対数になり，$\frac1{x^2}$ は指数が $-2$ なので，べきの公式を使う。答えを微分すると
\[
 \frac43\cdot\frac32x^{\frac12}
 +3x^{-2}=2\sqrt{x}+\frac3{x^2}
\]
となり，係数と符号の両方を確認できる。

和を項ごとに積分できるのは，微分が和と定数倍を保つためである。一方，積について
\[
 \int f(x)g(x)\dd x
 =\left(\int f(x)\dd x\right)
  \left(\int g(x)\dd x\right)
\]
という規則はない。右辺を微分すると積の微分の二項が現れ，ふつうは $f(x)g(x)$ にはならない。積があるときは，まず展開できるか，恒等式で和にできるかを考える。その後に学ぶ置換積分と部分積分も，この問題への別の答えである。

\section{一次式が入っているときの係数}\label{節:一次式の積分}

$\sin(3x)$ を積分する場合，$-\cos(3x)$ を候補にすると
\[
 \frac{\dd}{\dd x}\{-\cos(3x)\}=3\sin(3x)
\]
となる。合成関数の微分によって，内側の $3x$ の微分 $3$ が余分に掛かる。したがって候補を $3$ で割ればよい。
\[
 \int\sin(3x)\dd x=-\frac13\cos(3x)+C.
\]
一般に $F'=f$，$a\ne0$ ならば，
\[
 \frac{\dd}{\dd x}\left\{\frac1aF(ax+b)\right\}
 =\frac1a f(ax+b)\,a=f(ax+b).
\]
よって
\[
 \int f(ax+b)\dd x=\frac1aF(ax+b)+C
\]
となる。$a$ で割るのは書式上の約束ではなく，微分したときに現れる $a$ を打ち消すためである。

\begin{example}{内側の微分を確かめる}
次の不定積分を求めよ。
\[
 \int\left(4\e^{2x-1}+\frac3{5x+2}-2\cos(3x)\right)\dd x.
\]
\end{example}

それぞれ，指数の $2x-1$，分母の $5x+2$，角の $3x$ の微分を確かめる。これらは $2,5,3$ であるから，
\[
\begin{aligned}
 &\int\left(4\e^{2x-1}+\frac3{5x+2}-2\cos(3x)\right)\dd x\\
 &\qquad=2\e^{2x-1}+\frac35\log|5x+2|
       -\frac23\sin(3x)+C.
\end{aligned}
\]
$5x+2$ が $0$ になる点は除き，その左右の各区間で考える。対数の項を微分すると
\[
 \frac35\cdot\frac5{5x+2}=\frac3{5x+2}
\]
である。分母の係数 $5$ を見落とすと，ここで元の式へ戻らない。

一次式の場合は内側の微分が定数なので，係数の修正だけで済んだ。内側が $x^2+1$ ならば，微分は $2x$ であり，単なる定数倍ではない。次章の置換積分では，この $2x$ が積分される式の中にあるかどうかが手がかりになる。

\section{展開と割り算で和を作る}\label{節:展開と割算}

整式の積は，展開すれば各項が $x$ のべきになる。式を短く保つことと，積分しやすくすることは同じではない。積の形では微分の逆が見えなくても，和にすると基本公式を一項ずつ使える。

\begin{example}{積を展開する理由}
$\int x(2x-1)^2\dd x$ を求めよ。
\end{example}

二乗を展開してから $x$ を掛けると，
\[
 x(2x-1)^2=x(4x^2-4x+1)=4x^3-4x^2+x.
\]
したがって
\[
 \int x(2x-1)^2\dd x
 =x^4-\frac43x^3+\frac12x^2+C.
\]
候補を微分すれば $4x^3-4x^2+x$ に戻り，これをくくり直すと元の式になる。高いべきでは展開する項数が多くなるため，別の方法が有利になる。しかしこの例のように，数項で済むならば，展開は十分に見通しのよい方法である。

分数式では，分子の次数が分母以上のとき，まず整式の割り算を考える。割り算は，分母に掛ける整式と余りとに分子を分ける操作である。そうすれば分数全体が，「整式」と「分子の次数が低い分数」の和になる。

\begin{example}{分子の次数を下げる}
$x\ne-1$ において，$\int\frac{x^3+2x^2+3}{x+1}\dd x$ を求めよ。
\end{example}

分子を分母で割ると
\[
 x^3+2x^2+3=(x+1)(x^2+x-1)+4
\]
となる。この等式は，右辺を展開して確かめられる。両辺を $x+1$ で割って
\[
 \frac{x^3+2x^2+3}{x+1}
 =x^2+x-1+\frac4{x+1}
\]
とすれば，すべて既知の形である。よって
\[
 \int\frac{x^3+2x^2+3}{x+1}\dd x
 =\frac{x^3}{3}+\frac{x^2}{2}-x+4\log|x+1|+C.
\]
分数のまま眺めていたときに難しかったのは，高い次数の分子と分母が一つにまとまっていたためである。割り算によってその大部分を整式へ移すと，最後には一次式の逆数だけが残った。

約分できる式でも，もとの定義域は確かめる。たとえば $\frac{x^2-1}{x-1}=x+1$ は $x\ne1$ で成り立つ。右辺が $x=1$ でも定義されるからといって，左辺の関数が自動的にその点で定義されたことにはならない。積分する区間にその点が含まれるなら，連続に延長した関数を扱うのか，もとの式を扱うのかを明らかにする必要がある。

\section{部分分数で分母を分ける}\label{節:部分分数}

割り算を終えても，分母が二つの一次式の積になっていることがある。ここで分数の分母を勝手に分配することはできない。たとえば $\frac1{AB}$ は $\frac1A+\frac1B$ ではない。しかし，分子に合う係数を選べば，より簡単な分数の和として表せる。この書き換えを\textbf{部分分数分解}という。

\begin{example}{異なる一次因子に分ける}
$\int\frac{5x+1}{(x-1)(x+2)}\dd x$ を求めよ。
\end{example}

一次式の逆数ならば対数で積分できる。そこで
\[
 \frac{5x+1}{(x-1)(x+2)}
 =\frac A{x-1}+\frac B{x+2}
\]
と表せないかを調べる。右辺を通分すると分子は $A(x+2)+B(x-1)$ なので，求める係数は
\[
 5x+1=A(x+2)+B(x-1)
       =(A+B)x+(2A-B)
\]
を満たせばよい。係数を比べて
\[
 A+B=5,\qquad 2A-B=1
\]
を解くと $A=2$，$B=3$ となる。したがって
\[
 \int\frac{5x+1}{(x-1)(x+2)}\dd x
 =2\log|x-1|+3\log|x+2|+C.
\]
元の関数は $x=-2,1$ で定義されないので，答えはそれらをまたがない各区間で読む。

係数を決める途中で，分母を払った恒等式に $x=1,-2$ を代入してもよい。$x=1$ なら $6=3A$，$x=-2$ なら $-9=-3B$ となり，同じ係数が得られる。これは，元の分数に分母を $0$ にする値を代入したのではない。係数を決めるための整式の恒等式を調べているのである。

分母に同じ一次因子が重なっている場合も，低いべきから分けると積分できる。たとえば $\frac1{(x+1)^2}$ 自体はべきの公式で扱えるが，分子が一次式ならば，その分子を $x+1$ と定数の和にする。

\begin{example}{重なった一次因子を持つ分母}
$\int\frac{2x+3}{(x+1)^2}\dd x$ を求めよ。
\end{example}

分母にある $x+1$ を分子にも作りたいので，$2x+3=2(x+1)+1$ と書く。すると
\[
 \frac{2x+3}{(x+1)^2}
 =\frac2{x+1}+\frac1{(x+1)^2}.
\]
第一項には対数，第二項にはべきの公式を使い，
\[
 \int\frac{2x+3}{(x+1)^2}\dd x
 =2\log|x+1|-\frac1{x+1}+C
\]
を得る。二乗の分母があるからといって，すべてが逆数の原始関数になるわけではない。分解後の各項を見て，どの公式を使うかを決めている。

一般に有理関数，すなわち整式どうしの商を積分するときは，まず割り算で分子の次数を下げ，次に分母を因数分解して分けることを考える。ただし，実数の範囲で一次式に分かれない二次式もある。その場合に役立つ形を次に調べる。

\section{平方完成で二次式の中心を見る}\label{節:平方完成と積分}

二次式 $x^2-4x+5$ は
\[
 x^2-4x+5=(x-2)^2+1
\]
と平方完成できる。$x$ そのものよりも，$x-2$ を一つのまとまりとして見ると，二次式の中にある構造が明確になる。積分でも，このまとまりとその微分が式の中にあるかを調べると，候補が見つかることがある。

\begin{example}{平方根の微分と見比べる}
$\int\frac{x-2}{\sqrt{x^2-4x+5}}\dd x$ を求めよ。
\end{example}

平方完成すると，積分される式は
\[
 \frac{x-2}{\sqrt{(x-2)^2+1}}
\]
となる。分母に平方根があるので，$\sqrt{(x-2)^2+1}$ を微分してみる。合成関数の微分から
\[
 \frac{\dd}{\dd x}\sqrt{(x-2)^2+1}
 =\frac1{2\sqrt{(x-2)^2+1}}\,2(x-2)
 =\frac{x-2}{\sqrt{(x-2)^2+1}}.
\]
これがちょうど元の式なので，
\[
 \int\frac{x-2}{\sqrt{x^2-4x+5}}\dd x
 =\sqrt{x^2-4x+5}+C.
\]
根号の中は常に $1$ 以上だから，この答えは実数全体で使える。候補が見えた後に微分して一致を確かめたので，これは推測だけで終わった解答ではない。

同じ根号でも，分子に $x-2$ がない $\int\frac1{\sqrt{(x-2)^2+1}}\dd x$ では，いまの微分結果と一致しない。その違いを無視して同じ答えにはできない。また
\[
 \frac1{x^2-4x+3}=\frac1{(x-2)^2-1}
\]
ならば，平方の差なので $(x-3)(x-1)$ に分解できる。一方，$(x-2)^2+1$ は実数の一次式の積に分解できない。平方完成はいつでも積分を終わらせる操作ではないが，その後に因数分解を使えるのか，置換を考えるべきなのかを判断する材料になる。

\section{三角恒等式で積とべきをほどく}\label{節:三角式の整理}

$\sin x$ と $\cos x$ の原始関数は分かっていても，$\sin^2x$ や $\cos^2x$ は，いまの基本公式にそのまま当てはめられない。そこで，二倍角の公式
\[
 \cos(2x)=1-2\sin^2x=2\cos^2x-1
\]
を解き直して
\[
 \sin^2x=\frac{1-\cos(2x)}2,\qquad
 \cos^2x=\frac{1+\cos(2x)}2
\]
とする。左辺は二乗だが，右辺は定数と余弦の和になった。ここまで変形すれば基本公式を使える。

\begin{example}{二乗と積を和へ直す}
次の積分を求めよ。
\[
 \text{(1)}\quad\int_0^{\frac\pi2}\sin^2x\dd x,
 \qquad
 \text{(2)}\quad\int\sin(3x)\cos x\dd x.
\]
\end{example}

(1) は二倍角の公式から
\[
\begin{aligned}
 \int_0^{\frac\pi2}\sin^2x\dd x
 &=\frac12\int_0^{\frac\pi2}\{1-\cos(2x)\}\dd x\\
 &=\left[\frac x2-\frac14\sin(2x)\right]_0^{\frac\pi2}
 =\frac\pi4.
\end{aligned}
\]
$\cos(2x)$ の積分に $\frac12$ が付くため，全体の係数が $\frac14$ になる。最初から定積分を計算しているので，原始関数に任意定数を付けても上下端で差を取ると消える。

(2) は加法定理を二つ並べて考える。
\[
 \sin(3x+x)=\sin(3x)\cos x+\cos(3x)\sin x,
\]
\[
 \sin(3x-x)=\sin(3x)\cos x-\cos(3x)\sin x.
\]
加えると右辺の第二項が打ち消し合うので，
\[
 \sin(3x)\cos x=\frac12\{\sin(4x)+\sin(2x)\}.
\]
よって
\[
 \int\sin(3x)\cos x\dd x
 =-\frac18\cos(4x)-\frac14\cos(2x)+C.
\]
積を和へ直す公式を忘れても，加法定理から必要な式を作れる。微分して $\frac12\sin(4x)+\frac12\sin(2x)$ に戻ることを確かめ，もう一度積に直せば確認も終わる。

ここまでの計算では，積分記号の中の式を同じ値の別の式に書き換えてきた。どの変形を使うかは，元の式に現れた手がかりで決まる。整式の小さな積なら展開し，分子の次数が高ければ割り算をし，一次因子の積が分母にあれば部分分数を考える。二次式は平方完成でまとまりを見やすくし，三角関数の二乗や積は恒等式で和へ直す。

ただし，式を長く展開すれば必ずよいわけではない。たとえば $x(x^2+1)^{20}$ は展開できるが，$x^2+1$ を一つの変数と考えれば計算を短くできる。その見方を正当化するのが，次章で学ぶ置換積分である。

\section*{章末問題}
\begin{enumerate}[label=\textbf{\thechapter.\arabic*},leftmargin=*]
 \item\label{問:基本べき} $x>0$ において，$\int\frac{3x^2-2x+4}{\sqrt{x}}\dd x$ を求め，微分して確かめよ。
 \item\label{問:基本割算} $\int\frac{x^3+1}{x+1}\dd x$ を求めよ。式を整理した後も，もとの関数で除かれる点を述べよ。
 \item\label{問:基本部分分数} $\int\frac{5x-1}{(x-1)(x+1)}\dd x$ を求めよ。部分分数の係数を決める等式も書け。
 \item\label{問:基本重因子} $\int\frac{3x+4}{(x+1)^2}\dd x$ を求めよ。
 \item\label{問:基本平方完成} 分母を平方完成してから，$\int\frac1{x^2+6x+8}\dd x$ を求めよ。
 \item\label{問:基本三角二乗} $\int_0^{\frac\pi2}\cos^2(2x)\dd x$ を求めよ。
 \item\label{問:基本三角積} 加法定理から必要な積を和へ直す式を導き，$\int\sin(2x)\sin(5x)\dd x$ を求めよ。
\end{enumerate}

\smallskip\noindent 解答は\pageref{解答:第4章}ページ。
